\section{Polar curvature}
\label{sec:psh-polar-invariance}

Bounded changes preserve more than logarithmic coefficients: they preserve the whole curvature measure on the common pole set. A local energy estimate puts the bounded difference in a Sobolev space, while pole sets admit smooth tests of arbitrarily small energy. These tests exclude every coefficient measure of its Hessian from the pole set. Polar curvature is closed even when the pole set is not closed. A local elliptic equation therefore recovers it from its scalar Hermitian trace.

\begin{lemma}\label{lem:bounded-measure-laplacian-energy}
Let $U\subset\mathbb R^d$ be open and let $w\in L^\infty_{\mathrm{loc}}(U)$. If $\Delta w$ is a locally finite signed Radon measure, then $w\in W^{1,2}_{\mathrm{loc}}(U)$.
\end{lemma}
\begin{proof}
Fix $0\leq\psi\leq1$ smooth and compactly supported on $U$, and a relatively compact open neighborhood $V$ of its support. Mollify $w$ inside $V$, writing $w_\varepsilon=w*\rho_\varepsilon$ and $\nu_\varepsilon=(\Delta w)*\rho_\varepsilon$. For small $\varepsilon$, integration by parts and Young's inequality give
\[
 \begin{split}
 \int\psi^2|\nabla w_\varepsilon|^2
 &=-\int\psi^2w_\varepsilon\nu_\varepsilon
       -2\int\psi w_\varepsilon\nabla\psi\cdot\nabla w_\varepsilon\\
 &\leq \|w\|_{L^\infty(V)}|\Delta w|(V)
       +\tfrac12\int\psi^2|\nabla w_\varepsilon|^2
       +2\int w_\varepsilon^2|\nabla\psi|^2.
 \end{split}
\]
The right side bounds the local gradient norm independently of $\varepsilon$. Taking $\psi=1$ on any prescribed smaller compact set, weak Sobolev compactness and convergence of mollifications in $L^2$ identify the weak limit with $w$. This proves the assertion.
\end{proof}

\begin{lemma}\label{lem:pole-small-energy-tests}\index{Pluripolar set!smooth energy tests}
Let $v$ be quasi-plurisubharmonic on a complex manifold, let $U$ be a coordinate ball with compact closure in a larger chart, and let $K\subset U\cap\{v=-\infty\}$ be compact. For every open $O\subset U$ containing $K$, there are real $\varphi_j\in C^\infty_c(O)$ with $0\leq\varphi_j\leq1$, equal to one near $K$, such that $\|\varphi_j\|_{W^{1,2}(U)}\to0$. The norm uses the Euclidean coordinates.
\end{lemma}
\begin{proof}
A smooth quadratic correction and subtraction of a constant give a negative plurisubharmonic function $h$ on $U$ with the same pole set as $v$. Put $h_M=\max(h,-M)$ for $M\geq1$ and choose $0\leq\psi\leq1$ in $C^\infty_c(O)$, equal to one near $K$. On a fixed compact neighborhood of $\operatorname{supp}\psi$, the positive measures $\Delta h_M$ have uniformly bounded mass: choose a nonnegative smooth compactly supported $\chi$ on $U$ equal to one there and use
\[
 \int\chi\,\Delta h_M=\int h_M\Delta\chi,
 \qquad |h_M|\leq|h|\in L^1_{\mathrm{loc}}(U).
\]
Apply the preceding integration-by-parts estimate to mollifications of $h_M$. Their absolute values are at most $M$, and their squares are at most $M$ times their absolute values. The uniform Laplacian mass and local $L^1$ bounds therefore give, after passing to the weak Sobolev limit,
\[
 \int\psi^2|\nabla h_M|^2\leq CM,
 \qquad
 \left\|\frac{-\psi h_M}{M}\right\|_{W^{1,2}(U)}^2\leq\frac{C}{M}.
\]
Here $C$ depends on the fixed cutoffs and the local $L^1$ norm of $h$, but not on $M$. The second inequality also uses $h_M^2\leq M|h|$ in the terms containing $\psi$ and $\nabla\psi$.

The Sobolev function $-\psi h_M/M$ lies between zero and one. It equals one on a neighborhood of $K$, since $h$ is upper semicontinuous and $\{h<-M\}$ is an open neighborhood of $K$. Extend it by zero and mollify with sufficiently small radius. Its support stays in $O$, it remains equal to one on a smaller neighborhood of $K$, and its Sobolev norm does not increase. Taking $M\to\infty$ gives the required smooth functions.
\end{proof}

Restrictions below are taken coefficientwise on the Radon measures of a current of order zero. They therefore commute with smooth changes of coordinates.

\Needspace{10\baselineskip}
\begin{proposition}\label{prop:bounded-type-polar-invariance}\index{Singularity type!polar curvature invariance}\index{Polar curvature!bounded changes}
Let $u,v$ be quasi-plurisubharmonic on a complex manifold and locally equivalent up to bounded constants. Their pole sets agree, say $S=\{u=-\infty\}=\{v=-\infty\}$, and
\[
 \mathbf1_S\ddc u=\mathbf1_S\ddc v.
\]
Consequently, for any smooth real $(1,1)$-forms $\theta,\beta$,
\[
 \mathbf1_S(\theta+\ddc u)=\mathbf1_S(\beta+\ddc v).
\]
\end{proposition}
\begin{proof}
The two local order inequalities give the same pole set. Off $S$ put $w=u-v$ and assign arbitrary bounded values on $S$. This defines the locally bounded distribution $u-v$, because $S$ has Lebesgue measure zero. The coefficient measures of $\ddc w$ are locally finite: after a smooth correction, each Hessian of $u$ and $v$ is a positive current of order zero. Its Euclidean trace makes $\Delta w$ a signed measure. Thus \cref{lem:bounded-measure-laplacian-energy} gives $w\in W^{1,2}_{\mathrm{loc}}$.

Each real and imaginary part of a coefficient measure of $\ddc w$ is consequently a continuous functional on $W^{1,2}_0$ on a smaller ball, since a second distributional derivative of $w$ is a first derivative of its $L^2$ gradient. Let $\nu$ be any one of these real signed measures, with Jordan parts $\nu^+,\nu^-$. Choose a compact $K$ in the intersection of $S$ with a positive Hahn set, so $\nu^-(K)=0$. For $\epsilon>0$, take an open neighborhood $O$ of $K$ in the ball with $\nu^-(O)<\epsilon$. The tests of \cref{lem:pole-small-energy-tests}, applied to $u$, satisfy
\[
 \nu^+(K)\leq\int\varphi_j\,\mathrm d\nu^+
       =\langle\nu,\varphi_j\rangle
          +\int\varphi_j\,\mathrm d\nu^-
       \leq|\langle\nu,\varphi_j\rangle|+\epsilon.
\]
Continuity on the Sobolev space makes the first term tend to zero. Letting $\epsilon\downarrow0$ and using inner regularity gives $\nu^+(S)=0$ locally. Applying the same argument to $-\nu$ gives $\nu^-(S)=0$. All coefficient measures of $\ddc w$ therefore vanish on $S$, proving the first identity. Smooth forms also assign zero mass to $S$, which proves the second.
\end{proof}

For a closed positive current, local potentials turn equality of scalar traces into an elliptic equation. If both currents are carried by a Lebesgue null set, the resulting smooth difference vanishes.

\Needspace{10\baselineskip}
\begin{lemma}\label{lem:closed-null-current-trace}\index{Closed positive current!trace on a null set}
Let $T,S$ be closed positive $(1,1)$-currents on a complex manifold of dimension $n\geq1$, both carried coefficientwise by a Borel set of Lebesgue measure zero. For any Hermitian form $\eta$,
\[
 T\wedge\eta^{n-1}=S\wedge\eta^{n-1}
 \quad\Longrightarrow\quad T=S.
\]
\end{lemma}
\begin{proof}
On a smaller coordinate ball, take plurisubharmonic local potentials $a,b$ for $T,S$, by \cite[Chapter III, Proposition 1.19, p.~135]{DemCADG}, with the normalization of this book. Put $h=a-b\in L^1_{\mathrm{loc}}$. The trace identity gives $\ddc h\wedge\eta^{n-1}=0$ in distributions. Dividing by the smooth positive volume form $\eta^n$ gives a scalar second-order operator whose principal matrix is a positive multiple of the inverse Hermitian matrix of $\eta$. It is uniformly elliptic on every smaller ball, with smooth coefficients. Local elliptic regularity therefore makes $h$ smooth. Thus $T-S=\ddc h$ is smooth and carried by a Lebesgue null set. Its coefficient measures must vanish, proving equality locally and hence globally.
\end{proof}

Proper pushforward commutes with coefficientwise restriction to the inverse image of a Borel set. This gives exact compatibility of polar curvature under modifications, even when its pole set is not closed.

\Needspace{12\baselineskip}
\begin{proposition}\label{prop:polar-curvature-pushforward}\index{Polar curvature!pushforward}
Let $\pi:Y\to X$ be a modification of connected compact complex manifolds, let $\theta$ be a smooth real $(1,1)$-form on $X$, and let $u\in\PSH(X,\theta)$. Put $S=\{u=-\infty\}$ and
\[
 T=\theta+\ddc u,\qquad
 T_Y=\pi^*\theta+\ddc(u\circ\pi).
\]
Then
\begin{equation}\label{eq:polar-curvature-pushforward}
 \pi_*\bigl(\mathbf1_{\pi^{-1}(S)}T_Y\bigr)=\mathbf1_ST.
\end{equation}
\end{proposition}
\begin{proof}
Pullback is genuine and admissible by \cref{prop:psh-descent-modification}. Change of variables on the isomorphism locus, whose complement has Lebesgue measure zero on both spaces, gives $\pi_*(u\circ\pi)=u$ as distributions and $\pi_*\pi^*\theta=\theta$ as currents. Proper pushforward commutes with $\ddc$, so $\pi_*T_Y=T$.

For any Borel $B\subset X$, the coefficient measures of $\pi_*T_Y$ are pushforwards of the coefficient measures of $T_Y$ multiplied by the smooth coefficients arising from pullback of a test form. Restriction therefore gives $\pi_*(\mathbf1_{\pi^{-1}(B)}T_Y)=\mathbf1_B\pi_*T_Y$. Apply this to the Borel pole set $S$. No closedness of its restricted current is used.
\end{proof}

The closedness of polar curvature follows from the bounded cutoff argument for complete pluripolar restrictions. The cutoff approaches the complement of the pole set pointwise, so the pole set need not be closed.

\Needspace{12\baselineskip}
\begin{proposition}\label{prop:polar-curvature-closedness}\index{Polar curvature!closedness}
Let $X$ be a complex manifold of dimension $n\geq1$, let $u$ be quasi-plurisubharmonic on $X$, and let $\theta$ be a smooth real $(1,1)$-form. Then
\[
 \operatorname{Pol}_\theta(u)
       :=\mathbf1_{\{u=-\infty\}}(\theta+\ddc u)
\]
is a closed positive $(1,1)$-current.
\end{proposition}
\begin{proof}
Locally, a smooth quadratic correction and subtraction of a constant give a negative plurisubharmonic function $a$ with pole set $S=\{u=-\infty\}$. Put $\Theta=\ddc a$. Since smooth forms have no coefficient mass on $S$, the polar current equals $\mathbf1_S\Theta$ and is positive. In dimension one it is closed by degree.

Suppose $n\geq2$. Choose a smooth convex nondecreasing function $\chi$ on $\mathbb R$ with $\chi(t)=0$ for $t\leq1/2$ and $\chi(1)=1$, and a smooth function $0\leq\vartheta\leq1$ which is zero near zero and one near one. Set
\[
 w_k=\chi(e^{a/k}),\qquad k\geq1.
\]
These functions are bounded plurisubharmonic, with $0\leq w_k\leq1$. They vanish on $S$ and tend to one at every point outside $S$. Hence $\vartheta(w_k)\Theta$ converges coefficientwise locally in total variation to $\mathbf1_{X\setminus S}\Theta$. Also $\vartheta'(w_k)$ vanishes on $S$ and tends to zero everywhere.

Fix a constant strongly positive $(n-2,n-2)$-form $\beta$ and put $\Xi=\Theta\wedge\beta$. The bounded Bedford--Taylor chain rule gives
\[
 \mathrm d\bigl(\vartheta(w_k)\Xi\bigr)
       =\vartheta'(w_k)\,\mathrm dw_k\wedge\Xi.
\]
Indeed, decrease smooth plurisubharmonic functions to $w_k$ on a smaller ball. Dominated convergence gives convergence of their $\vartheta$-weighted currents. The bounded Bedford--Taylor energy convergence and Cauchy--Schwarz give convergence of the differentiated expressions, proving the chain rule for $w_k$.
For a smooth compactly supported test one-form $\alpha$, choose a nonnegative smooth cutoff $\tau$ equal to one near its support. Since $w_k\geq0$ and $\Xi$ is closed and positive, bounded integration by parts gives
\[
 2\int\tau\,\mathrm dw_k\wedge\dc w_k\wedge\Xi
       \leq\int\tau\,\ddc(w_k^2)\wedge\Xi
       =\int w_k^2\ddc\tau\wedge\Xi\leq C.
\]
Here $C$ depends only on the local mass of $\Xi$ and the fixed cutoff. The Cauchy--Schwarz inequality for a positive current now yields
\[
 \bigl|\langle\vartheta'(w_k)\,\mathrm dw_k\wedge\Xi,
                    \alpha\rangle\bigr|^2
       \leq C_\alpha\int_{\operatorname{supp}\alpha}
                 |\vartheta'(w_k)|^2\,\Xi\wedge\kappa
       \longrightarrow0,
\]
where $\kappa$ is the Euclidean Hermitian form. The limit follows by dominated convergence against the fixed locally finite trace measure. Thus $\mathbf1_{X\setminus S}\Theta\wedge\beta$ is closed. Constant strongly positive forms span the needed test forms, so $\mathbf1_{X\setminus S}\Theta$ is closed. Subtracting it from $\Theta$ proves closedness of $\mathbf1_S\Theta$. The local currents agree on overlaps and give the asserted global current.
\end{proof}
